\textbf{VALUTARE SE METTERE GRAFICI}


\section{Graph Concepts}
A \textbf{Graph} is a collection of \textbf{nodes} (vertices) and \textbf{edges}.
\begin{definition}[Adjacent nodes]
    Two nodes are adjacent if there is an edge between them    
\end{definition}
\begin{definition}[Complete graph]
    A graph is complete if there is an edge between every pair of nodes   
\end{definition}
\begin{definition}[Directed graph]
    A graph is directed\footnote{orientato} if the all his edges are directed   
\end{definition}
In a directed graph we can see some sets of nodes, based on the relationships showed by the edges:
\begin{itemize}
    \item \textbf{pa(X)} is the parent set of X
    \item \textbf{ch(Y)} is the child set of Y
    \item \textbf{de(Y)} is the descendants of Y, that is the set of nodes which can be reached by Y
    \item \textbf{an(Z)} is the set of ancestors of Z, the set of nodes from which Z can be reached.
\end{itemize}
It is also important to distinguish between a \textbf{Path}, that is any sequence of adjacent nodes, 
regardless of the direction of the edges that join them, and a \textbf{Directed Path}, that is a path
between two nodes traced along the arrows. So as a conseguence, in a directed path it can't be any node
on the path that has two edges on the path directed into it or out of it.
\newline
\section{From PGMs to CGMs}
The \textit{Probabilitstic Grapical models} are statistical models more often represented by 
\textbf{Bayesian Networks} in which we want to study the associations or the correlations between 
variables by modelling the \textbf{Data Distribution}:
\begin{definition}[Chain Rule]
    \begin{equation}
        \begin{aligned}
            P(X_1, X_2, \dots, X_n) &= P(X_1) P(X_2 \mid X_1) P(X_3 \mid X_1, X_2) \cdots P(X_n \mid X_1, \dots, X_{n-1}) \\
            &= P(X_1) \prod_{i=2}^{n} P(X_i \mid X_1, \dots, X_{i-1})
        \end{aligned}
    \end{equation}
\end{definition}
